\section{CONTINUOUS GROUPS AND INFINITESIMAL TRANSFORMATIONS.}

Suddenly, in the autumn of 1873, Lie takes up again the study of transformation groups and obtains decisive results. For as much as it is possible to follow the unravelling of his thoughts in a few letters to A. Mayer from the years 1873-1874 ({[}202{]}, vol.~5, pp.~585-608), he starts from a ``continuous group'' of transformations on n variables

\[
x _ {i} ^ {\prime} = f _ {i} (x _ {1}, \dots , x _ {n}, a _ {1}, \dots , a _ {r}) (1 \leq i \leq n)\tag{25.4}
\]

depending effectively \(^{3}\) on r parameters \(a_{1},\ldots,a_{r}\) ; he observes that if the transformation (25.4) is the identity for the values \(a_{1}^{0},\ldots,a_{r}^{0}\) of the parameters, \(^{4}\) then the Taylor expansions of the \(x_{i}\) , limited to the first order:

\[
f _ {i} \left(x _ {1}, \dots , x _ {n}, a _ {1} ^ {0} + z _ {1}, \dots , a _ {r} ^ {0} + z _ {r}\right) = x _ {i} + \sum_ {k = 1} ^ {r} z _ {k} X _ {k i} \left(x _ {1}, \dots , x _ {n}\right) + \dots (1 \leq i \leq n)\tag{25.5}
\]

give a ``generic'' infinitely small transformation depending linearly on the r parameters \(z_{j}\)

\[
d x _ {i} = \left(\sum_ {k = 1} ^ {r} z _ {k} X _ {k i} (x _ {1}, \dots , x _ {n})\right) d t (1 \leq i \leq n).\tag{25.6}
\]

Proceeding as in his memoir with Klein, Lie integrates the differential system

\[
\frac {d \xi_ {1}}{\sum_ {k} z _ {k} X _ {k 1} (\xi_ {1} , \dots , \xi_ {n})} = \dots = \frac {d \xi_ {n}}{\sum_ {k} z _ {k} X _ {k n} (\xi_ {1} , \dots , \xi_ {n})} = d t,\tag{25.7}
\]

which gives him, for each point \((z_{1},\ldots,z_{r})\) , a group with one parameter

\[
t \mapsto x _ {i} ^ {\prime} = g _ {i} (x _ {1}, \dots , x _ {n}, z _ {1}, \dots , z _ {r}, t) (1 \leq i \leq n)\tag{25.8}
\]

such that \(g_{i}(x_{1},\ldots,x_{n},z_{1},\ldots,z_{r},0)=x_{i}\) for every i. He shows in an ingenious way, using the fact that the transformations (25.4) make up a set closed under composition, that the group with one parameter (25.8) is a subgroup of the given group ({[}202{]}, vol.~5, Abh. VII, pp.~32-63). The new idea, key to the whole theory, is to proceed on to the second order in the Taylor expansions of the functions (25.4). The thread of his argument is fairly confused and heuristic ({[}202{]}, vol.~5, Abh. VII, pp.~32-63 and {[}202{]}, vol.~5, pp.~600-601); it can be set out as follows. For the \(z_{j}\) that are fairly small, one can put t=1 in (25.8), and thus are obtained new parameters \(z_{1},\ldots,z_{r}\) for the transformations of the group (it is in fact the first appearance of the ``canonical parameters''). By definition, we have with regard to (25.7)

\[
\frac {\partial g _ {i}}{\partial t} = \sum_ {k} z _ {k} X _ {k i} (x _ {1} ^ {\prime}, \dots , x _ {n} ^ {\prime}),
\]

from where

\[
\frac {\partial^ {2} g _ {i}}{\partial t ^ {2}} = \sum_ {k, j} z _ {k} \frac {\partial X _ {k i}}{\partial x _ {j}} (x _ {1} ^ {\prime}, \dots , x _ {n} ^ {\prime}) \frac {\partial x _ {j} ^ {\prime}}{\partial t} =
\]

\[
\sum_ {k, j} z _ {k} \frac {\partial X _ {k i}}{\partial x _ {j}} (x _ {1} ^ {\prime}, \dots , x _ {n} ^ {\prime}) \left(\sum_ {h} z _ {h} X _ {h j} (x _ {1} ^ {\prime}, \dots , x _ {n} ^ {\prime})\right)
\]

which gives

\[
\begin{array}{l} x _ {i} ^ {\prime} = x _ {i} + \left(\sum_ {k} z _ {k} X _ {k i} (x _ {1}, \dots , x _ {n})\right) t \\ \quad + \frac {1}{2} \left(\sum_ {k, h, j} z _ {k} z _ {h} \frac {\partial X _ {k i}}{\partial x _ {j}} (x _ {1}, \dots , x _ {n}) X _ {h j} (x _ {1}, \dots , x _ {n})\right) t ^ {2} + \dots , \end{array}
\]

from where, for \(t = 1\) , the Taylor expansions with respect to the parameters \(z_{j}\)

\[
x _ {i} ^ {\prime} = x _ {i} + \left(\sum_ {k} z _ {k} X _ {k i}\right) + \frac {1}{2} \left(\sum_ {k, h, j} z _ {k} z _ {h} X _ {h j} \frac {\partial X _ {k i}}{\partial x _ {j}}\right) + \dots (1 \leq i \leq n).\tag{25.9}
\]

Let us write these relations in brief as \(x' = G(x, z)\) between vectors

\[
\boldsymbol {x} = (x _ {1}, \dots , x _ {n}), \boldsymbol {x} ^ {\prime} = (x _ {1} ^ {\prime}, \dots , x _ {n} ^ {\prime}), z = (z _ {1}, \dots , z _ {r});
\]

the fundamental property of closure of the sets of these transformations under composition is written as

\[
G (G (\boldsymbol {x}, \boldsymbol {u}), \boldsymbol {v}) = G (\boldsymbol {x}, H (\boldsymbol {u}, \boldsymbol {v}))\tag{25.10}
\]

where \(H = (H_{1},\ldots ,H_{r})\) is independent of \(x\) ; it is immediate that \(H(u,0) = u,H(0,v) = v\) , whence the expansions

\[
H _ {i} (u, v) = u _ {i} + v _ {i} + \frac {1}{2} \sum_ {h, k} c _ {i k h} u _ {h} v _ {k} + \dots ,\tag{25.11}
\]

the terms not written down being non linear either in u or in v. Transforming (25.10) with the help of (25.9) and (25.11), then comparing the terms in \(u_{h}v_{k}\) of the two members, Lie obtains the relations

\[
\sum_ {j = 1} ^ {n} \left(X _ {h j} \frac {\partial X _ {k i}}{\partial x _ {j}} - X _ {k j} \frac {\partial X _ {h i}}{\partial x _ {j}}\right) = \sum_ {l = 1} ^ {r} c _ {l h k} X _ {l i} (1 \leq h, k \leq r, 1 \leq i \leq n).\tag{25.12}
\]

His expertise with the theory of partial differential equations leads him to write these conditions in a simpler form: following the model of (25.1), he associates with each of these r infinitely small transformations obtained by putting \(z_{k}=1\) , \(z_{h}=0\) for \(h\neq k\) in (25.6), the differential operator

\[
A _ {k} (f) = \sum_ {i = 1} ^ {n} X _ {k i} \frac {\partial f}{\partial x _ {i}},\tag{25.13}
\]

and rewrites the conditions (25.12) in the form

\[
[ A _ {h}, A _ {k} ] = \sum_ {l} c _ {l h k} A _ {l},\tag{25.14}
\]

a keystone of his theory. Until then, he had used indiscriminately the terms ``infinitely small transformation'' and ``infinitesimal transformation'' (e. g. {[}202{]}, vol.~5, Abh. I, pp.~1-4); the simplicity of the relations (25.14) leads him to call the operator (25.13) the ``symbol'' of the infinitesimal transformation \(dx_{i} = X_{ki}dt\) ( \(1 \leq i \leq n\) ) ({[}202{]}, vol.~5, Abh. III, pp.~42-75) and very rapidly, it is the operator (25.13) itself that he will call ``infinitesimal transformation'' ({[}202{]}, vol.~5, Abh. III, pp.~42-75 and {[}202{]}, vol.~5, p.~589).

He then becomes aware of the tight links which unite the theory of ``continuous groups'' with his previous research on contact transformations and partial differential equations. This link fills him with enthusiasm: ``My old work was so to speak all ready in advance to found the new theory of groups of transformations'' he wrote to Mayer in 1874 ({[}202{]}, vol.~5, p.~586).

In the following years, Lie pursued the study of transformation groups. Apart from the general theorems summarised here (§ III), he obtains a certain number of more particular results: the determination of transformation groups of the straight line and of the plane, subgroups of small codimension in projective groups, groups with at most 6 parameters, etc.. For all that he does not abandon differential equations. In fact, it even appears that for him, the theory of transformation groups was to be an instrument for integrating differential equations, where the transformation group would play a role analogous to that of the Galois group of an algebraic equation. \(^{5}\) Let us note that this research leads him likewise to introducing certain sets of transformations with an infinity of parameters, which he calls ``infinite and continuous groups''; \(^{6}\) he keeps the name ``finite and continuous groups'' for the transformation groups with a finite number of parameters of type (25.4) above.

